% GENERATED FILE. Edit canonical lessons, then run npm run generate:books.
% canonical-source-hash: fnv1a64:d1102fc993c2b1d0
% canonical-lessons: PA-C247-ala, PA-C247-bari, PA-C247-khunja, PA-C247-pirha, PA-C247-sanduk

\chapter{Inside the Room}
\label{ch:pa-a2-inside-the-room}
\hlchaptermodality{247}

\begin{chapteropening}
\textbf{By the end of this chapter:} \emph{I can name places to keep things in a room.}

Complete the last lesson of chapter 247: I can name places to keep things in a room.
\end{chapteropening}

\section[\texorpdfstring{\pa{ਆਲਾ} (ālā)}{ਆਲਾ (ālā)}]{\pa{ਆਲਾ} (ālā) — a wall niche}

\label{lesson:PA-C247-ala}



\begin{quote}
Before the new one: say the Punjabi for a roof gutter, then the Punjabi for a high shelf.
\end{quote}

\subsection*{You'll want to know: \pa{ਆਲਾ}}
\textbf{\pa{ਆਲਾ}} — \emph{ālā} — \textquotedblleft{}a wall niche\textquotedblright{}.

\begin{quote}
\pa{ਆਲੇ} \pa{ਵਿੱਚ} \pa{ਦੀਵਾ} \pa{ਹੈ।}
\end{quote}

\begin{quote}
\emph{āle vich dīvā hai.}
\end{quote}

\begin{quote}
There is a lamp in the wall niche.
\end{quote}

A masculine noun.

\begin{cousinweb}
\textbf{\pa{ਆਲਾ}} (\emph{ālā}) is inherited from Sanskrit \emph{ālaya}, \textquotedblleft{}house, receptacle\textquotedblright{}; CDIAL lists Punjabi \emph{ālā}, \textquotedblleft{}niche in a wall\textquotedblright{} (CDIAL 1366).
\end{cousinweb}

\begin{grammarlens}[title={What you've built}]
One more word for this chapter.
\end{grammarlens}

\subsection*{Your turn}
\begin{itemize}
\raggedright
  \item \emph{Say it:} \emph{ālā}
  \item \emph{Say it:} \emph{ālā}, once more
  \item \emph{Say it:} \emph{āle vich dīvā hai.} — There is a lamp in the wall niche.
  \item \emph{Recall:} say the Punjabi for a roof gutter, then the Punjabi for a high shelf, then say \emph{ālā} again
\end{itemize}

\begin{tcolorbox}[breakable,colback=teal!4,colframe=teal!35!black,title={Before you move on}]
What does \pa{ਆਲਾ} mean? (\textquotedblleft{}A wall niche\textquotedblright{}.) Say it once more.
\end{tcolorbox}
\section[\texorpdfstring{\pa{ਬਾਰੀ} (bārī)}{ਬਾਰੀ (bārī)}]{\pa{ਬਾਰੀ} (bārī) — a small window}

\label{lesson:PA-C247-bari}



\begin{quote}
Before the new one: say the Punjabi for a high shelf, then the Punjabi for a wall niche.
\end{quote}

\subsection*{You'll want to know: \pa{ਬਾਰੀ}}
\textbf{\pa{ਬਾਰੀ}} — \emph{bārī} — \textquotedblleft{}a small window\textquotedblright{}.

\begin{quote}
\pa{ਰਸੋਈ} \pa{ਦੀ} \pa{ਬਾਰੀ} \pa{ਛੋਟੀ} \pa{ਹੈ।}
\end{quote}

\begin{quote}
\emph{rasoī dī bārī chhoṭī hai.}
\end{quote}

\begin{quote}
The kitchen window is small.
\end{quote}

A feminine noun.

\begin{cousinweb}
\textbf{\pa{ਬਾਰੀ}} (\emph{bārī}) is inherited from Sanskrit \emph{dvāra}, \textquotedblleft{}door\textquotedblright{}; CDIAL lists Punjabi \emph{bārī}, \textquotedblleft{}window\textquotedblright{} (CDIAL 6663).
\end{cousinweb}

\begin{grammarlens}[title={What you've built}]
One more word for this chapter.
\end{grammarlens}

\subsection*{Your turn}
\begin{itemize}
\raggedright
  \item \emph{Say it:} \emph{bārī}
  \item \emph{Say it:} \emph{bārī}, once more
  \item \emph{Say it:} \emph{rasoī dī bārī chhoṭī hai.} — The kitchen window is small.
  \item \emph{Recall:} say the Punjabi for a high shelf, then the Punjabi for a wall niche, then say \emph{bārī} again
  \item \emph{Write it:} copy \textbf{\pa{ਬਾਰੀ}} once, checking the whole word
\end{itemize}

\begin{tcolorbox}[breakable,colback=teal!4,colframe=teal!35!black,title={Before you move on}]
What does \pa{ਬਾਰੀ} mean? (\textquotedblleft{}A small window\textquotedblright{}.) Say it once more.
\end{tcolorbox}
\section[\texorpdfstring{\pa{ਖੂੰਜਾ} (khūnjā)}{ਖੂੰਜਾ (khūnjā)}]{\pa{ਖੂੰਜਾ} (khūnjā) — a corner of a room}

\label{lesson:PA-C247-khunja}



\begin{quote}
Before the new one: say the Punjabi for a wall niche, then the Punjabi for a small window.
\end{quote}

\subsection*{You'll want to know: \pa{ਖੂੰਜਾ}}
\textbf{\pa{ਖੂੰਜਾ}} — \emph{khūnjā} — \textquotedblleft{}a corner of a room\textquotedblright{}.

\begin{quote}
\pa{ਝਾੜੂ} \pa{ਕਮਰੇ} \pa{ਦੇ} \pa{ਖੂੰਜੇ} \pa{ਵਿੱਚ} \pa{ਹੈ।}
\end{quote}

\begin{quote}
\emph{jhāṛū kamre de khūnje vich hai.}
\end{quote}

\begin{quote}
The broom is in the corner of the room.
\end{quote}

A masculine noun.

\begin{cousinweb}
\textbf{\pa{ਖੂੰਜਾ}} (\emph{khūnjā}) is inherited from a reconstructed \emph{*khōñcā}, \textquotedblleft{}curve, bend\textquotedblright{}; CDIAL lists Punjabi \emph{khū̃jā}, \textquotedblleft{}corner\textquotedblright{} (CDIAL 3930).
\end{cousinweb}

\begin{grammarlens}[title={What you've built}]
One more word for this chapter.
\end{grammarlens}

\subsection*{Your turn}
\begin{itemize}
\raggedright
  \item \emph{Say it:} \emph{khūnjā}
  \item \emph{Say it:} \emph{khūnjā}, once more
  \item \emph{Say it:} \emph{jhāṛū kamre de khūnje vich hai.} — The broom is in the corner of the room.
  \item \emph{Recall:} say the Punjabi for a wall niche, then the Punjabi for a small window, then say \emph{khūnjā} again
  \item \emph{Write it:} copy \textbf{\pa{ਖੂੰਜਾ}} once, checking the whole word
\end{itemize}

\begin{tcolorbox}[breakable,colback=teal!4,colframe=teal!35!black,title={Before you move on}]
What does \pa{ਖੂੰਜਾ} mean? (\textquotedblleft{}A corner of a room\textquotedblright{}.) Say it once more.
\end{tcolorbox}
\section[\texorpdfstring{\pa{ਪੀੜ੍ਹਾ} (pīṛhā)}{ਪੀੜ੍ਹਾ (pīṛhā)}]{\pa{ਪੀੜ੍ਹਾ} (pīṛhā) — a low stool}

\label{lesson:PA-C247-pirha}



\begin{quote}
Before the new one: say the Punjabi for a small window, then the Punjabi for a corner of a room.
\end{quote}

\subsection*{You'll want to know: \pa{ਪੀੜ੍ਹਾ}}
\textbf{\pa{ਪੀੜ੍ਹਾ}} — \emph{pīṛhā} — \textquotedblleft{}a low stool\textquotedblright{}.

\begin{quote}
\pa{ਦਾਦੀ} \pa{ਪੀੜ੍ਹੇ} \pa{ਉੱਤੇ} \pa{ਬੈਠਦੀ} \pa{ਹੈ।}
\end{quote}

\begin{quote}
\emph{dādī pīṛhe utte baiṭhdī hai.}
\end{quote}

\begin{quote}
Grandmother sits on a low stool.
\end{quote}

A masculine noun.

\begin{cousinweb}
\textbf{\pa{ਪੀੜ੍ਹਾ}} (\emph{pīṛhā}) is inherited from Sanskrit \emph{pīṭha}, \textquotedblleft{}stool, bench\textquotedblright{}; CDIAL lists Punjabi \emph{pīṛhā}, \textquotedblleft{}stool\textquotedblright{} (CDIAL 8222).
\end{cousinweb}

\begin{grammarlens}[title={What you've built}]
One more word for this chapter.
\end{grammarlens}

\subsection*{Your turn}
\begin{itemize}
\raggedright
  \item \emph{Say it:} \emph{pīṛhā}
  \item \emph{Say it:} \emph{pīṛhā}, once more
  \item \emph{Say it:} \emph{dādī pīṛhe utte baiṭhdī hai.} — Grandmother sits on a low stool.
  \item \emph{Recall:} say the Punjabi for a small window, then the Punjabi for a corner of a room, then say \emph{pīṛhā} again
  \item \emph{Write it:} copy \textbf{\pa{ਪੀੜ੍ਹਾ}} once, checking the whole word
\end{itemize}

\begin{tcolorbox}[breakable,colback=teal!4,colframe=teal!35!black,title={Before you move on}]
What does \pa{ਪੀੜ੍ਹਾ} mean? (\textquotedblleft{}A low stool\textquotedblright{}.) Say it once more.
\end{tcolorbox}
\section[\texorpdfstring{\pa{ਸੰਦੂਕ} (sandūk)}{ਸੰਦੂਕ (sandūk)}]{\pa{ਸੰਦੂਕ} (sandūk) — a trunk; a storage chest}

\label{lesson:PA-C247-sanduk}



\begin{quote}
Before the new one: say the Punjabi for a corner of a room, then the Punjabi for a low stool.
\end{quote}

\subsection*{You'll want to know: \pa{ਸੰਦੂਕ}}
\textbf{\pa{ਸੰਦੂਕ}} — \emph{sandūk} — \textquotedblleft{}a trunk; a storage chest\textquotedblright{}.

\begin{quote}
\pa{ਸੰਦੂਕ} \pa{ਵਿੱਚ} \pa{ਰਜਾਈਆਂ} \pa{ਹਨ।}
\end{quote}

\begin{quote}
\emph{sandūk vich rajāīāṁ han.}
\end{quote}

\begin{quote}
There are quilts in the trunk.
\end{quote}

A masculine noun.

\begin{cousinweb}
\textbf{\pa{ਸੰਦੂਕ}} (\emph{sandūk}) is borrowed: it is usually traced to Persian \emph{sandūq}, \textquotedblleft{}a chest, a box\textquotedblright{}.
\end{cousinweb}

\begin{grammarlens}[title={What you've built}]
One more word for this chapter.
\end{grammarlens}

\subsection*{Your turn}
\begin{itemize}
\raggedright
  \item \emph{Say it:} \emph{sandūk}
  \item \emph{Say it:} \emph{sandūk}, once more
  \item \emph{Say it:} \emph{sandūk vich rajāīāṁ han.} — There are quilts in the trunk.
  \item \emph{Recall:} say the Punjabi for a corner of a room, then the Punjabi for a low stool, then say \emph{sandūk} again
  \item \emph{Write it:} copy \textbf{\pa{ਸੰਦੂਕ}} once, checking the whole word
\end{itemize}

\begin{tcolorbox}[breakable,colback=teal!4,colframe=teal!35!black,title={Before you move on}]
What does \pa{ਸੰਦੂਕ} mean? (\textquotedblleft{}A trunk; a storage chest\textquotedblright{}.) Say it once more.
\end{tcolorbox}
